\section{Discussão}

Os resultados mostram um claro compromisso entre acurácia preditiva e restrições de implantação. No benchmark com múltiplas sementes, Random Forest e LSTM alcançaram menores erros macro médios, enquanto a SNN permaneceu competitiva e claramente mais estável que o MLP. \rev{Os resultados de Wilcoxon e Friedman, junto com os intervalos de confiança sobrepostos, indicam que Random Forest, LSTM e SNN são estatisticamente indistinguíveis neste conjunto de dados, o que reforça que a SNN deve ser interpretada como uma alternativa robusta orientada à borda, e não como um modelo de melhor acurácia universal.} Do ponto de vista metodológico, a configuração multimodal mostrou-se valiosa porque mantém alvos fisicamente relacionados em um único pipeline de inferência e simplifica a lógica de decisão subsequente em aplicações de campo.

Da perspectiva de sistemas, o comportamento das plataformas é heterogêneo: o Dispositivo B alcançou a maior vazão, o Dispositivo C obteve a menor energia incremental por inferência e o ESP32 forneceu viabilidade offline completa a taxas de processamento muito menores. Como a potência de base das plataformas varia com a arquitetura e o comportamento do sistema operacional, a energia por inferência combinada com a vazão é um critério mais acionável do que watts absolutos isoladamente. A ablação do aumento de dados mostrou melhora modesta mas consistente do MAE macro (2,53\%) e ganho em turbidez, indicando que a expansão sintética controlada ajuda a estabilizar o treinamento neste regime de poucas amostras.

\subsection{Diretrizes Práticas de Implantação}
Os resultados combinados permitem uma diretriz simples de implantação para equipes de campo: para amostragem temporal densa, atualizações frequentes de mapas ou painéis quase em tempo real, a execução em smartphones ou SBCs é preferível; para verificações periódicas esparsas com estrita simplicidade de hardware, o ESP32 permanece viável e operacionalmente atraente, apesar dos maiores valores de mJ/inf.

Um segundo ponto prático é a política de seleção de modelos: em cenários de campo restritos, selecionar apenas pelo MAE mínimo pode ser subótimo. Um critério melhor é um envelope multiobjetivo que combine erro de predição aceitável, comportamento estável entre sementes, vazão da plataforma e energia incremental por inferência, extensível com restrições como orçamento de bateria e cadência de telemetria.

\rev{A escolha da SNN em vez do Random Forest merece justificativa explícita, já que o Random Forest alcançou o menor MAE médio. Primeiro, a diferença de acurácia não é estatisticamente significativa: os testes de Wilcoxon e os intervalos de confiança mostram sobreposição completa entre os dois modelos. Segundo, a SNN oferece propriedades que o Random Forest não possui: computação orientada a eventos cujo custo escala com a atividade de pulsos, um único modelo denso compacto que exporta diretamente para ONNX/TFLite em todas as três classes de plataforma, e compatibilidade futura com aceleradores neuromórficos, nos quais a execução pulsada esparsa pode reduzir a energia bem abaixo das medições em CPU aqui reportadas. Terceiro, o parâmetro de codificação temporal $T$ fornece um controle ajustável de acurácia--energia sem contrapartida em comitês de árvores. A SNN é, portanto, selecionada não como o modelo de menor erro, mas como o modelo estatisticamente equivalente com o melhor envelope de implantação.}

\rev{Quanto a Edge AI versus Cloud AI, o compromisso vai além da computação bruta. A inferência em nuvem oferece modelos maiores e atualizações mais simples, mas em implantações remotas na Amazônia adiciona a energia de rádio, que tipicamente domina o orçamento por amostra, latência e disponibilidade atreladas a enlaces intermitentes, custos recorrentes de conectividade e exposição de dados ambientais brutos a infraestrutura de terceiros. A inferência de borda totalmente offline inverte esse perfil: a energia é limitada pela computação local (aqui medida na escala de $\mu$J a mJ por inferência), o monitoramento continua durante interrupções e apenas resultados agregados precisam de sincronização eventual. Um padrão híbrido --- triagem local com consolidação oportunista em nuvem --- é uma evolução natural do presente pipeline.}

\subsection{Interpretação de Erros e Validade}
Os gráficos de predito vs.\ real indicam um claro perfil de erro dependente do alvo: pH e OD seguem a tendência de identidade com menor dispersão, enquanto a turbidez exibe espalhamento mais amplo. Esse comportamento é coerente com a heterogeneidade de águas doces e reforça que a turbidez deve ser interpretada com margens de confiança mais amplas no uso operacional. Da perspectiva de aplicação, essa assimetria apoia uma estratégia de decisão em camadas: tratar pH e OD como sinais de triagem de curto prazo mais fortes, e usar a turbidez primariamente para sinalizar zonas candidatas a amostragem adicional, e não como gatilho isolado. \rev{Para o oxigênio dissolvido especificamente, a análise por alvo mostra que nenhum modelo avaliado alcança $R^2$ médio positivo, o que aponta para uma limitação do conjunto de atributos, e não da modelagem: as entradas atuais carregam pouca informação relacionada ao OD (\mbox{$|r| \leq 0,24$}). Melhorar a predição de OD exigirá, portanto, enriquecer a modalidade de entrada --- p.~ex., condutividade, potencial de oxirredução, clorofila ou perfis diários de temperatura --- em vez de mais ajuste de modelo.}

As ameaças à validade permanecem centradas no tamanho amostral e na representatividade regional. Embora o aumento de dados, a avaliação com múltiplas sementes e o benchmarking multiplataforma tenham melhorado a robustez, a generalização mais ampla ainda requer campanhas em múltiplas estações e múltiplos sítios, com protocolos de aquisição harmonizados e validação externa repetida. \rev{A transferibilidade para outras regiões deve ser enquadrada em nível de pipeline, e não dos pesos treinados: os gradientes espaciais aprendidos são específicos do Lago do Serpa, de modo que a aplicação direta do modelo ajustado a outro lago não deve preservar a acurácia. O que se transfere é a metodologia --- projeto de atributos, política de aumento, protocolo estatístico e pilha de implantação ---, que pode ser retreinada localmente a partir de uma campanha comparavelmente pequena em um novo sítio. Estratégias de ajuste fino entre sítios e de adaptação de domínio estão planejadas para quantificar explicitamente esse custo de transferência.} A reprodutibilidade é uma força prática deste pipeline: o protocolo de modelagem, a avaliação baseada em sementes e as medições de implantação offline são explícitos e portáveis, permitindo replicação direta em outros contextos de baixa conectividade. Além das métricas preditivas, a implantação no Raspberry Pi alcançou inferência submilissegundo e a validação no ESP32 confirmou execução offline autônoma sob recursos restritos, reforçando que a inferência local é praticável em perfis de borda heterogêneos.
